\documentclass[9pt]{extarticle}
\usepackage[twocolumn,margin=0.55in,columnsep=0.28in]{geometry}
\usepackage{amsmath,amssymb}
\usepackage{booktabs}
\usepackage{microtype}
\usepackage[dvipsnames]{xcolor}
\usepackage[colorlinks=true,linkcolor=NavyBlue,citecolor=NavyBlue,urlcolor=NavyBlue]{hyperref}
\usepackage{enumitem}
\setlist{nosep,leftmargin=1.1em,itemsep=0pt,topsep=1pt}
\pagestyle{empty}
\setlength{\parskip}{2pt}

% Generated by paper/make_numbers.py -- do not edit by hand.
% Every measurement quoted in the paper is defined here from
% results/data/*.json, so the prose cannot drift from the data.

\newcommand{\EfiveAbsReffMax}{4.6}
\newcommand{\EfiveAbsReffMin}{2.3}
\newcommand{\EfiveDarkAUC}{0.561}
\newcommand{\EfiveDarkEff}{98.4}
\newcommand{\EfiveDomEff}{17.4}
\newcommand{\EfiveFisherAUC}{0.775}
\newcommand{\EfiveFisherEff}{9.2}
\newcommand{\EfiveKLgbar}{0.301}
\newcommand{\EfiveKLprobe}{0.019}
\newcommand{\EfiveMeanReff}{0.31}
\newcommand{\EfiveModels}{gpt2, distilgpt2, pythia-70m}
\newcommand{\EfiveN}{400}
\newcommand{\EfiveNfamilies}{2}
\newcommand{\EfiveNmodels}{3}
\newcommand{\EfiveSecond}{distilgpt2}
\newcommand{\EfiveSecondDarkEff}{98.8}
\newcommand{\EfiveSecondDomEff}{13.7}
\newcommand{\EfiveSecondFisherEff}{7.3}
\newcommand{\EfiveThirdDarkEff}{98.6}
\newcommand{\EfiveThirdDomEff}{14.1}
\newcommand{\EfiveThirdFisherEff}{8.2}
\newcommand{\EfiveThirdName}{pythia-70m}
\newcommand{\EfiveWidths}{512, 768}
\newcommand{\EfourAucNinety}{14.7}
\newcommand{\EfourAucNinetyList}{32/6/6}
\newcommand{\EfourAucNinetyNine}{15.3}
\newcommand{\EfourCertNinety}{1.3}
\newcommand{\EfourCertNinetyList}{2/1/1}
\newcommand{\EfourCertNinetyNine}{1.3}
\newcommand{\EfourLayers}{3, 6, 9}
\newcommand{\EfourM}{64}
\newcommand{\EfourNlayers}{3}
\newcommand{\EfourOracleNinety}{1.0}
\newcommand{\EfourOracleNinetyList}{1/1/1}
\newcommand{\EfourOracleNinetyNine}{1.0}
\newcommand{\EfourRandNinety}{29.3}
\newcommand{\EfourRandNinetyList}{46/22/20}
\newcommand{\EfourRandNinetyNine}{49.0}
\newcommand{\EfourSpAUC}{0.313}
\newcommand{\EfourSpSbar}{0.761}
\newcommand{\EfourSpeedupAUC}{11}
\newcommand{\EfourSpeedupRandom}{22}
\newcommand{\EoneAxisPopEffect}{$0$}
\newcommand{\EoneCorrP}{0.061}
\newcommand{\EoneCorrS}{0.124}
\newcommand{\EoneDarkAUC}{0.502}
\newcommand{\EoneDarkGap}{$5\times10^{-14}$}
\newcommand{\EoneEmpEffect}{$2.0\times10^{-2}$}
\newcommand{\EonePopEffect}{$2\times10^{-15}$}
\newcommand{\EoneProbeAUC}{0.998}
\newcommand{\EoneProbeAUCclosed}{0.998}
\newcommand{\EsixLateMatched}{0.062}
\newcommand{\EsixLateUnmatched}{0.039}
\newcommand{\EsixMatched}{0.031}
\newcommand{\EsixN}{300}
\newcommand{\EsixNbeh}{5}
\newcommand{\EsixNlab}{4}
\newcommand{\EsixNlayers}{12}
\newcommand{\EsixRatio}{0.97}
\newcommand{\EsixUnmatched}{0.032}
\newcommand{\EsixWins}{6}
\newcommand{\EthreeD}{768}
\newcommand{\EthreeDarkAUC}{0.676}
\newcommand{\EthreeDialLevels}{0.50, 0.70, 0.85, 0.95}
\newcommand{\EthreeDialRhos}{0.066, 0.051, 0.057, 0.047}
\newcommand{\EthreeEffDark}{99.7}
\newcommand{\EthreeEffProbe}{6.0}
\newcommand{\EthreeEffRandom}{2.6}
\newcommand{\EthreeFPRidentZero}{50}
\newcommand{\EthreeFPRneutZero}{0}
\newcommand{\EthreeGbarAUC}{0.698}
\newcommand{\EthreeHostileGap}{15.5}
\newcommand{\EthreeIdentAUClast}{0.520}
\newcommand{\EthreeInWarranty}{10}
\newcommand{\EthreeMaxAUC}{1.000}
\newcommand{\EthreeMeanReff}{0.18}
\newcommand{\EthreeMeanRho}{0.056}
\newcommand{\EthreeMinRho}{0.006}
\newcommand{\EthreeModel}{gpt2}
\newcommand{\EthreeNlayers}{12}
\newcommand{\EthreePairGap}{1.044}
\newcommand{\EthreePairN}{96}
\newcommand{\EthreePairSign}{71.9}
\newcommand{\EthreeSpAUC}{-0.063}
\newcommand{\EthreeSpSbar}{0.917}
\newcommand{\EtwoAbsReff}{1.6}
\newcommand{\EtwoAcc}{1.000}
\newcommand{\EtwoD}{64}
\newcommand{\EtwoDarkHi}{98.8}
\newcommand{\EtwoDarkLo}{98.7}
\newcommand{\EtwoDeltaInSD}{0.2}
\newcommand{\EtwoEffHi}{12.0}
\newcommand{\EtwoEffList}{12.0, 12.1, 12.0, 12.0}
\newcommand{\EtwoEffLo}{12.0}
\newcommand{\EtwoLayers}{5}
\newcommand{\EtwoLevels}{0.50, 0.70, 0.90, 0.98}
\newcommand{\EtwoMarkerList}{0.504, 0.703, 0.899, 0.979}
\newcommand{\EtwoNlevels}{4}
\newcommand{\EtwoNruns}{12}
\newcommand{\EtwoNseeds}{3}
\newcommand{\EtwoRhoDelta}{-0.004}
\newcommand{\EtwoRhoHi}{0.147}
\newcommand{\EtwoRhoList}{0.151, 0.150, 0.148, 0.147}
\newcommand{\EtwoRhoLo}{0.151}
\newcommand{\EtwoSeedSD}{0.017}
\newcommand{\PfiveMargin}{0.081}
\newcommand{\PfiveNrand}{400{,}000}
\newcommand{\PsevenAltRtwo}{0.721}
\newcommand{\PsevenAltSlope}{0.904}
\newcommand{\PsevenRtwo}{0.959}
\newcommand{\PsevenSlope}{1.006}
\newcommand{\PthreeIso}{0.975}
\newcommand{\PthreeMaxRelErr}{$1.3\times10^{-2}$}

\newcommand{\Sbar}{\bar{S}}
\newcommand{\gbar}{\bar{g}}

\title{\vspace{-2.4em}\textbf{\large Seeing Is Not Steering}\\[-0.1em]
{\normalsize A one-page case for effect-based steering certificates}\vspace{-1.6em}}
\author{\normalsize Blessings Mambwe \quad \texttt{bleymambwe@gmail.com}\vspace{-2.6em}}
\date{}

\begin{document}
\twocolumn[
  \maketitle
  \vspace{-1.2em}
]

\paragraph{The problem, without jargon.} To check whether a language model
``represents'' a concept, practitioners fit a linear classifier (a
\emph{probe}) to its hidden states. To control a behaviour, they add a
vector to the hidden state and watch the output change (\emph{activation
steering}). Standard practice reuses the probe's direction as the steering
vector. These are two different jobs --- separating labelled examples, and
moving an output --- and nothing in how a probe is fit ties them together.
No widely used method checks whether they coincide before either is
deployed.

\paragraph{Precise problem statement.} Fix a scalar behaviour $\varphi$
(here, a logit difference tracking toxic output) and a hidden state
$x_\ell$ at layer $\ell$. For a unit-norm direction $w$ and step $\alpha$,
the causal effect of writing $\alpha w$ into $x_\ell$ is
$\Delta(\alpha,w)=\varphi(x_\ell+\alpha w)-\varphi(x_\ell)$. At fixed write
norm, how much of the best achievable $\Delta$ do four candidate
directions actually deliver --- \textbf{random}, the \textbf{probe}
direction, a direction built \textbf{orthogonal} to the probe, and the
\textbf{best} direction found by directly measuring many candidates --- and
can the answer be predicted \emph{before} any of them is tried?

\paragraph{One cheap number, zero interventions.} Let $\gbar$ be the average
gradient of $\varphi$ with respect to $x_\ell$, from one backward pass over
the same data the probe already required. Define the \textbf{certificate}
$\Sbar(w) := |\gbar^\top w|/\|w\|$. Computing $\Sbar$ needs no intervention
at all.

\paragraph{Headline result.} On GPT-2, a toxicity probe that separates
held-out prompts at AUC \EthreeMaxAUC{} recovers only \EthreeEffProbe\% of
the best available causal effect; a probe-orthogonal direction recovers
\EthreeEffDark\%; a random direction recovers \EthreeEffRandom\%
(Table~\ref{tab:main}). The pattern replicates on RealToxicityPrompts
across three models and two architecture families, scoring the exact
difference-of-means vector activation-steering methods deploy
(\EfiveDomEff\%, versus \EfiveDarkEff\% for the orthogonal direction). The
certificate rank-predicts measured effect (Spearman \EthreeSpSbar) where
probe accuracy does not (\EthreeSpAUC); ordering candidate interventions by
it needs $\sim\EfourSpeedupRandom\times$ fewer measurements than random
search to find the best available direction.

\begin{table}[t]
\centering
\small
\begin{tabular}{lrr}
\toprule
Direction & GPT-2 (\%) & RealToxicityPrompts (\%) \\
\midrule
Random & \EthreeEffRandom & --- \\
Probe & \EthreeEffProbe & \EfiveFisherEff \\
Difference-of-means & --- & \EfiveDomEff \\
Probe-orthogonal & \EthreeEffDark & \EfiveDarkEff \\
Certificate-optimal & 100 & 100 \\
\bottomrule
\end{tabular}
\caption{\% of best available causal effect recovered at equal write norm.
Certificate-optimal is the normalising basis, not a measured percentage.}
\label{tab:main}
\end{table}

\paragraph{Why the gap exists.} The probe direction is fit to separate
label-$1$ from label-$0$ examples; it depends only on the class means and
covariance of the hidden states. The certificate-optimal direction depends
only on how $\varphi$ responds to a perturbation. Nothing forces these two
disjoint quantities to point the same way, and on the systems measured here
they very nearly do not.

\paragraph{Systematic experiments and de-risking.}
\begin{itemize}
\item De-risk on an exactly solvable linear synthetic system with known
ground truth, before touching a real model.
\item Validate on GPT-2 with matched sentence pairs that isolate the causal
factor (hostile framing) from a spurious confound (identity terms), so the
effect is not a dataset artefact.
\item Replicate on three models, two architecture families, on a public
toxicity benchmark, scoring the exact vectors deployed methods ship.
\item Rule out ``this is just a bad probe'': retrain it on an unrelated
concept and check whether alignment changes --- it does not
($\rho=\EsixMatched$ matched vs.\ $\rho=\EsixUnmatched$ unmatched, ratio
\EsixRatio).
\item Test whether $\Sbar$ can replace trial-and-error search: order
candidates by $\Sbar$ under a fixed measurement budget and compare to
random and probe-accuracy ordering.
\end{itemize}

\paragraph{Industry relevance.} Many deployed safety monitors are built
from exactly this kind of probe. If probe accuracy does not predict
steering efficacy, probe-based monitoring is not, by itself, evidence that
a model's behaviour is under control. $\Sbar$ gives a one-backward-pass
number that predicts which of many candidate directions will actually
work, in principle replacing the expensive trial-and-error search teams
currently run by hand.

\paragraph{Related work.} Activation steering and its deployed
difference-of-means vectors~\cite{actadd,caa}; local-linearity optimal
control for steering~\cite{locallin2026}; identifiability limits of
circuit-level mechanistic interpretability~\cite{eeao2025}; control-theoretic
(Gramian-based) framings of neural interpretability~\cite{blackbox2025}.
Full positioning against 12 further references in the companion report.

\paragraph{Full paper.} Formal definitions, an exact decoupling theorem, a
closed form for probe-orthogonal directions, a finite-sample leakage law,
and complete experimental and reproducibility detail are in the companion
technical report, \emph{Seeing Is Not Steering} (Mambwe, 2026). Every
number above is generated from \texttt{results/data/*.json}; none is typed
by hand.

\begin{thebibliography}{9}
\small
\bibitem{actadd}
A.~M.~Turner, L.~Thiergart, G.~Leech, D.~Udell, U.~Mini, and M.~MacDiarmid.
\emph{Activation Addition: Steering Language Models Without Optimization.}
2024. arXiv:2308.10248.
\bibitem{caa}
N.~Rimsky, N.~Gabrieli, J.~Schulz, M.~Tong, E.~Hubinger, and A.~M.~Turner.
\emph{Steering Llama 2 via Contrastive Activation Addition.} ACL, 2024.
arXiv:2312.06681.
\bibitem{locallin2026}
J.~Skifstad, X.~A.~Yang, and G.~Chou. \emph{Local Linearity of LLMs Enables
Activation Steering via Model-Based Linear Optimal Control.} 2026.
arXiv:2604.19018.
\bibitem{eeao2025}
M.~Méloux, F.~Portet, S.~Maniu, and M.~Peyrard. \emph{Everything, Everywhere,
All at Once: Is Mechanistic Interpretability Identifiable?} ICLR, 2025.
arXiv:2502.20914.
\bibitem{blackbox2025}
J.~Moon. \emph{From Black-Box to White-Box: Control-Theoretic Neural Network
Interpretability.} 2025. arXiv:2511.12852.
\end{thebibliography}

\end{document}
